% Stage 8 restored fragment — 2.4.3 Uniqueness of radiative solutions (archive lines 5456-5772)
\subsection{Uniqueness of radiative solutions}
\label{sec:ch243}

Outgoing selection is not a post hoc spatial mask: Sommerfeld radiation and uniqueness on $\mathcal{B}_\Gamma$ decide which Maxwell solutions carry watt-bearing export before harmonic charts attach \cite{sommerfeld1949,stratton1941}. Reactive near-field energy may be locally large while contributing negligibly to $P_{\mathrm{rad}}$ on $\Gamma_R$ once $\Pi_{\mathrm{rad}}$ is fixed \cite{harrington2001}.
Without exterior uniqueness, two distinct field pairs could satisfy the same feed data while
differing on the measurement sphere---a failure no wattmeter or pattern probe could resolve
\cite{sommerfeld1949}. The proof below is therefore not a mathematical luxury
imposed for elegance; it is the physical condition under which laboratory attribution of
radiation to declared sources becomes meaningful. Outgoing selection is not a post hoc spatial
mask: Sommerfeld radiation and uniqueness on
$\mathcal{B}_\Gamma$ decide which Maxwell solutions carry watt-bearing export before harmonic
charts attach \cite{sommerfeld1949}. On $\Gamma_R$, reactive near-field energy may be locally
large while contributing negligibly to $P_{\mathrm{rad}}$ once $\mathcal{R}_\omega$ and
$\Pi_{\mathrm{rad}}$ are fixed \cite{harrington2001}. Exterior uniqueness ties measured patterns
to declared sources within the radiative class \cite{stratton1941}. Solvers that omit outgoing
filters fit non-radiating currents that compact ranges cannot distinguish from export without
watt completion on $\mathcal{R}_\omega[\Jimp]$ rather than on $\mathcal{S}_\omega[\Jimp]$ alone
\cite{jackson1999}. The proof below imports the difference-field energy method from
Subsection~\ref{sec:ch01.1.3} and specializes it to the Sommerfeld outgoing class fixed in
Subsection~\ref{sec:ch241}.

Subsection~\ref{sec:ch243} proves that exterior radiative solutions are unique on the
outgoing branch, upgrading hypothesis
Eq.~\eqref{eq:ch2.exterior-uniqueness-hypothesis} to exterior radiative uniqueness
\cite{stratton1941}. The difference-field reduction
Eqs.~\eqref{eq:ch1.uniqueness-difference-system}--\eqref{eq:ch1.uniqueness-null} and
energy identity Eq.~\eqref{eq:ch1.uniqueness-energy-identity} from
Subsection~\ref{sec:ch01.1.3} are presupposed; the passive uniqueness argument and
coercive estimate Eq.~\eqref{eq:ch1.uniqueness-estimate} are not re-proved.
Subsection~\ref{sec:ch243} specializes the energy method to exterior domains with
Eqs.~\eqref{eq:ch2.sommerfeld-field-condition}--\eqref{eq:ch2.sommerfeld-outgoing-class},
outgoing resolvent selection
Eqs.~\eqref{eq:ch2.outgoing-resolvent-selection}--\eqref{eq:ch2.outgoing-solution-operator},
and radiation refinement Eq.~\eqref{eq:ch2.outgoing-radiation-refinement}. Source
reconstruction limits belong to Subsection~\ref{sec:ch244}. Let \(\Omega_{\mathrm{ext}}\) denote an exterior domain with compact
\(\operatorname{supp}\Jimp\subset B_R(0)\), vacuum constants \(\varepsilon_0,\mu_0\)
in the far zone, and fixed \(\omega>0\) in the harmonic band of
Subsection~\ref{sec:ch223}. Data are admissible per
Eq.~\eqref{eq:ch1.admissible-data-class} with outgoing radiation condition
Eq.~\eqref{eq:ch2.sommerfeld-outgoing-class} at infinity, equivalently
Eq.~\eqref{eq:ch1.radiation-trace-condition} on spherical shells
\cite{stratton1941}. Constitutive response in the source region is passive:
\(\Im\{\varepsilon\}\ge 0\), \(\Im\{\mu\}\ge 0\) or equivalent loss so homogeneous
difference energy is dissipative \cite{stratton1941}. Incoming and advanced branches
are excluded by Eq.~\eqref{eq:ch2.outgoing-selector-def} and
Eq.~\eqref{eq:ch2.incoming-outgoing-complement} \cite{harrington2001}. Fields lie in
\(\mathcal{F}_{\mathrm{SF}}\) as defined in Subsection~\ref{sec:ch241}; radiative
export uses \(\mathcal{F}_{\mathrm{rad}}\subset\mathcal{F}_{\mathrm{SF}}\) from
Eq.~\eqref{eq:ch2.far-zone-transverse-radiation} \cite{stratton1941}.
Uniqueness is established by showing that two outgoing solutions for the same impressed data
cannot differ nontrivially: the proof therefore studies the homogeneous problem satisfied by
their difference \cite{stratton1941}. Let \((\E_1,\H_1)\) and \((\E_2,\H_2)\) be two harmonic
solutions on the outgoing branch for the same impressed data and \(\mathcal{B}_\Gamma\). Define
\(\Delta\E=\E_1-\E_2\), \(\Delta\H=\H_1-\H_2\), and
\(\mathbf{u}_\Delta=(\Delta\E,\Delta\H)\). Linearity and
Eq.~\eqref{eq:ch1.uniqueness-difference-system} give the homogeneous exterior system
\begin{equation}
\label{eq:ch2.exterior-difference-system}
\curl\Delta\E=-\jj\omega\mu\,\Delta\H,\quad
\curl\Delta\H=\jj\omega\varepsilon\,\Delta\E,\quad
\mathcal{B}_\Gamma[\Delta\E,\Delta\H]=\mathbf{0},
\end{equation}
with \(\varepsilon,\mu\) matching the physical model in \(\Omega_s\) and
\(\varepsilon_0,\mu_0\) outside \cite{stratton1941}. Physically, if two outgoing solutions
for the same feed existed, their difference would satisfy this homogeneous system while still
carrying export on $\Gamma_R$---a situation the flux argument below excludes
\cite{stratton1941}. Outgoing radiation condition requires
\begin{equation}
\label{eq:ch2.exterior-difference-outgoing}
\mathcal{O}_{\mathrm{out}}[\mathbf{u}_\Delta]=\mathbf{u}_\Delta,
\qquad
\mathbf{u}_\Delta\in\mathcal{F}_{\mathrm{SF}},
\end{equation}
which is the exterior specialization of the third row of
Eq.~\eqref{eq:ch1.uniqueness-difference-system} when
\(\mathcal{B}_\Gamma\) includes Sommerfeld radiation \cite{harrington2001}. Equations
\eqref{eq:ch2.exterior-difference-system}--\eqref{eq:ch2.exterior-difference-outgoing}
state that uniqueness testing reduces to showing only the trivial pair satisfies the
homogeneous outgoing problem \cite{stratton1941}. Incoming difference modes would
violate Eq.~\eqref{eq:ch2.exterior-difference-outgoing} even if they satisfied the
curl rows \cite{jackson1999}.
On truncated exterior shells, form the complex Poynting vector for difference fields,
\(\mathbf{S}_\Delta=\tfrac12\Delta\E\times\Delta\H^{*}\), as in
Eq.~\eqref{eq:ch1.uniqueness-energy-identity} \cite{stratton1941}. For a large
sphere \(\Gamma_R=\{r=R\}\) enclosing \(\Omega_s\), let
\(\Omega_R=B_R\setminus\Omega_s\). The divergence theorem on \(\Omega_R\) yields
\begin{equation}
\label{eq:ch2.exterior-flux-identity}
\int_{\Gamma_R}\hat{\mathbf{r}}\cdot\mathbf{S}_\Delta\,dS
=
-\jj\omega\!\int_{\Omega_R}
\Big(u_{\Delta H}-u_{\Delta E}\Big)\,dV
\;+\;\mathcal{P}_{\Delta,s},
\end{equation}
where \(u_{\Delta E}=\tfrac14\varepsilon|\Delta\E|^{2}\),
\(u_{\Delta H}=\tfrac14\mu|\Delta\H|^{2}\), and \(\mathcal{P}_{\Delta,s}\) collects
source-region passive dissipation when \(\Im\{\varepsilon\},\Im\{\mu\}>0\)
\cite{stratton1941}. Equation~\eqref{eq:ch2.exterior-flux-identity}
is Eq.~\eqref{eq:ch1.uniqueness-energy-identity} with \(\Delta\Jimp=\mathbf{0}\) and
\(\Gamma=\Gamma_R\); its real part controls uniqueness through the sign of the
boundary flux \cite{stratton1941}. On a compact range, the real part of the $\Gamma_R$ integral
is what a wattmeter would read from the difference field if such a superposition were
physically present; outgoing selection forces that reading to nonnegative export or zero
\cite{harrington2001}. In lossless vacuum on
\(\Omega_R\setminus\Omega_s\), the volume term is purely imaginary at steady state,
so \(\Re\int_{\Gamma_R}\hat{\mathbf{r}}\cdot\mathbf{S}_\Delta\,dS\) is determined
entirely by far-zone radiation \cite{harrington2001}.
Sommerfeld condition Eq.~\eqref{eq:ch2.sommerfeld-field-condition} implies that on
\(\Gamma_R\) with \(R\gg a=\max|\mathbf{r}|\) over \(\operatorname{supp}\Jimp\), the
leading transverse contribution to \(\Delta\E,\Delta\H\) scales as
\(e^{-\jj k_0 R}/R\) with \(k_0=\omega\sqrt{\varepsilon_0\mu_0}\)
\cite{stratton1941}. The far-zone transverse radiation law
Eq.~\eqref{eq:ch2.far-zone-transverse-radiation} then gives
\begin{equation}
\label{eq:ch2.exterior-flux-asymptotic}
\Re\int_{\Gamma_R}\hat{\mathbf{r}}\cdot\mathbf{S}_\Delta\,dS
=
P_{\Delta,\mathrm{rad}}(R)
=
\frac{1}{2\eta_0}\oint_{\Gamma_R}
\big|\mathbf{E}_{\Delta,t}\big|^{2}\,dS
\;+\;\mathcal{O}(R^{-1}),
\end{equation}
with \(\eta_0=\sqrt{\mu_0/\varepsilon_0}\) and \(\mathbf{E}_{\Delta,t}\) the
transverse spherical component of \(\Delta\E\) \cite{stratton1941}. Taking
\(R\to\infty\) under outgoing radiation condition,
\begin{equation}
\label{eq:ch2.exterior-sommerfeld-flux-limit}
\lim_{R\to\infty}P_{\Delta,\mathrm{rad}}(R)
=
P_{\Delta,\mathrm{rad}}\ge 0,
\end{equation}
where \(P_{\Delta,\mathrm{rad}}\) is the net outgoing power carried by the difference
field \cite{stratton1941}. Equation
\eqref{eq:ch2.exterior-sommerfeld-flux-limit} is the exterior boundary term induced
by Eq.~\eqref{eq:ch1.radiation-trace-condition}: outgoing difference fields radiate
nonnegative power to infinity, while incoming difference fields would produce negative
flux and are excluded by \(\mathcal{O}_{\mathrm{out}}\)
\cite{jackson1999}. Dyadic asymptotics
Eq.~\eqref{eq:ch2.sommerfeld-dyadic-asymptotic} are presupposed; they are not
re-tabulated here \cite{harrington2001}.
\begin{figure}[t]
\centering
\ChOneFigBlock{%
\begin{tikzpicture}[ch1fig, node distance=5mm and 7mm]
 \node[ch1box] (s) {Compact\\$\Omega_s$};
 \node[ch1box, right=16mm of s] (d) {$\Delta\E,\Delta\H$\\outgoing};
 \node[ch1box, right=16mm of d] (gr) {$\Gamma_R$\\flux};
 \node[ch1box, right=16mm of gr] (out) {$P_{\Delta,\mathrm{rad}}\ge 0$};
 \draw[ch1flow] (s) -- (d) -- (gr) -- (out);
\end{tikzpicture}%
}
\caption{Exterior uniqueness flux argument: homogeneous difference fields on the
outgoing branch satisfy Eqs.~\eqref{eq:ch2.exterior-flux-identity}--\eqref{eq:ch2.exterior-sommerfeld-flux-limit}
on spheres \(\Gamma_R\).}
\label{fig:ch2.exterior-uniqueness-flux}
\end{figure}
Section~\ref{sec:ch01.3}: volume terms in
Eq.~\eqref{eq:ch2.exterior-flux-identity} are controlled by passivity when loss is
present, and vanish in lossless vacuum away from sources \cite{stratton1941}. \phantomsection\label{lem:ch2.zero-flux-difference}
Let \(\mathbf{u}_\Delta=(\Delta\E,\Delta\H)\) satisfy
Eqs.~\eqref{eq:ch2.exterior-difference-system}--\eqref{eq:ch2.exterior-difference-outgoing}
on an exterior domain with vacuum far zone. If
\(P_{\Delta,\mathrm{rad}}=0\) in Eq.~\eqref{eq:ch2.exterior-sommerfeld-flux-limit},
then \(\mathbf{E}_{\Delta,t}=\mathbf{0}\) on every sphere \(\Gamma_R\) for
\(R\) large enough, and \(\mathbf{u}_\Delta\) has no propagating outgoing content in
\(\mathcal{F}_{\mathrm{rad}}\) \cite{stratton1941}.
Eq.~\eqref{eq:ch2.exterior-flux-asymptotic} with \(P_{\Delta,\mathrm{rad}}=0\) forces
\(\oint_{\Gamma_R}|\mathbf{E}_{\Delta,t}|^{2}\,dS=0\) for all sufficiently large
\(R\). Outgoing transverse fields with zero spherical flux cannot carry
\(1/R\) radiation; the only remaining homogeneous outgoing transverse modes are
identically zero on \(\mathcal{F}_{\mathrm{rad}}\) \cite{stratton1941}. \phantomsection\label{thm:ch2.exterior-radiative-uniqueness}
Let \(\Omega_{\mathrm{ext}}\) be an exterior domain with compact impressed sources,
vacuum far zone, and outgoing Sommerfeld completion
Eqs.~\eqref{eq:ch2.sommerfeld-outgoing-class}--\eqref{eq:ch2.outgoing-selector-def}.
If \((\E_1,\H_1)\) and \((\E_2,\H_2)\) are harmonic Maxwell solutions on the outgoing
branch for the same admissible data, then \(\E_1=\E_2\) and \(\H_1=\H_2\).
Equivalently,
\begin{equation}
\label{eq:ch2.exterior-uniqueness-theorem}
\mathcal{L}_\omega[\mathbf{u}_\omega]=\mathbf{0}
\ \text{and}\ 
\mathbf{u}_\omega\in\mathcal{F}_{\mathrm{SF}}
\ \Longrightarrow\
\mathbf{u}_\omega=\mathbf{0},
\end{equation}
which upgrades Eq.~\eqref{eq:ch2.exterior-uniqueness-hypothesis} to the Chapter~\ref{ch:02}
definitive uniqueness statement \cite{stratton1941}.
Let \(\mathbf{u}_\Delta\) be the difference of two outgoing solutions. Then
Eqs.~\eqref{eq:ch2.exterior-difference-system}--\eqref{eq:ch2.exterior-difference-outgoing}
hold. Apply Eq.~\eqref{eq:ch2.exterior-flux-identity} and take
\(\Re\{\cdot\}\). In lossless vacuum outside \(\Omega_s\), the volume integral is
purely imaginary, so \(\Re P_{\Delta,\mathrm{rad}}(R)\) equals the real boundary flux.
Outgoing radiation condition forces \(\Re P_{\Delta,\mathrm{rad}}\ge 0\); passivity in
\(\Omega_s\) adds nonnegative \(\Re\mathcal{P}_{\Delta,s}\), so the only way to
satisfy Eq.~\eqref{eq:ch2.exterior-sommerfeld-flux-limit} with equality throughout is
\(P_{\Delta,\mathrm{rad}}=0\) \cite{stratton1941}. The auxiliary relation (\ref{lem:ch2.zero-flux-difference})
eliminates outgoing transverse content. Evanescent components remain in
\(\mathcal{F}_{\mathrm{SF}}\) but are tied to the same source-free outgoing solution
operator; with no radiative export and homogeneous forcing, the unique retarded
exterior state is trivial by Eq.~\eqref{eq:ch2.outgoing-solution-operator} with
\(\Jimp=\mathbf{0}\) \cite{harrington2001}. Hence
\(\mathbf{u}_\Delta=\mathbf{0}\), proving Eq.~\eqref{eq:ch2.exterior-uniqueness-theorem}.
Exterior radiative uniqueness (Subsection~\ref{sec:ch243}) is the exterior counterpart of
the passive uniqueness argument (Subsection~\ref{sec:ch01.1.3}): instead of a bounded passive
\(\Gamma\) with coercivity Eq.~\eqref{eq:ch1.uniqueness-estimate}, the exterior uses
Sommerfeld flux Eq.~\eqref{eq:ch2.exterior-sommerfeld-flux-limit} as the coercive
boundary channel \cite{stratton1941}. The uniqueness index
Eq.~\eqref{eq:ch1.uniqueness-condition-index} is finite on exterior outgoing problems
away from embedded resonances in finite scatterers; it diverges on bounded cavities at
Eq.~\eqref{eq:ch2.cavity-resonance-failure} \cite{harrington2001}. Exterior uniqueness
therefore replaces bounded coercivity with an \emph{open} boundary condition that
selects decaying outgoing phase \cite{stratton1941}.
Exterior radiative uniqueness (Subsection~\ref{sec:ch243}) makes the physical resolvent
injective on the outgoing subspace:
\begin{equation}
\label{eq:ch2.exterior-resolvent-injective}
\mathcal{L}_\omega^{-1}\big|_{\mathrm{phys}}:
\mathcal{Y}_\omega\to\mathcal{F}_{\mathrm{SF}}
\ \text{is single-valued},
\qquad
\mathcal{S}_\omega[\Jimp]
=
\mathcal{L}_\omega^{-1}\big|_{\mathrm{phys}}
\begin{bmatrix}\mathbf{s}_\omega(\Jimp,\rho)\\ \mathbf{g}_\Gamma\end{bmatrix},
\end{equation}
consistent with Eq.~\eqref{eq:ch2.green-bvp-equivalence}, outgoing Green branch
Eq.~\eqref{eq:ch2.outgoing-green-branch}, and solution selection
Eq.~\eqref{eq:ch2.outgoing-solution-operator} \cite{stratton1941}.
\phantomsection\label{cor:ch2.outgoing-green-unique}
For fixed \(\Jimp\) in the admissible class, the exterior field
\(\mathcal{S}_\omega[\Jimp]\) obtained from
Eq.~\eqref{eq:ch2.outgoing-solution-operator} is unique in
\(\mathcal{F}_{\mathrm{SF}}\) \cite{stratton1941}.
Two convolutions differing by a homogeneous outgoing state violate
exterior radiative uniqueness (Subsection~\ref{sec:ch243}) \cite{harrington2001}.
Subsection~\ref{sec:ch201} declared \(\Pi_{\mathrm{rad}}\) idempotent on
\(\mathcal{F}_{\mathrm{rad}}\); the joint property with outgoing radiation condition can now be
stated precisely. On radiative states,
\begin{equation}
\label{eq:ch2.rad-idempotency-radiative}
\Pi_{\mathrm{rad}}^2=\Pi_{\mathrm{rad}}
\quad\text{on}\ \mathcal{F}_{\mathrm{rad}},
\end{equation}
as developed in Subsection~\ref{sec:ch201} \cite{stratton1941}. After
exterior radiative uniqueness (Subsection~\ref{sec:ch243}), the measurement sequence
Eq.~\eqref{eq:ch2.outgoing-radiation-refinement} defines a single-valued map
\(\mathcal{R}_\omega=\Pi_{\mathrm{rad}}\circ\mathcal{S}_\omega\) without branch
ambiguity from duplicate outgoing decompositions \cite{harrington2001}. Evanescent
components remain in \(\ker\Pi_{\mathrm{rad}}\) through
Eq.~\eqref{eq:ch2.rad-null-on-evanescent}; uniqueness of \(\mathcal{S}_\omega\) does
not imply uniqueness of \(\mathcal{R}_\omega[\Jimp]\) when near-zone evanescent
storage dominates export \cite{stratton1941}.
Outgoing selector idempotency Eq.~\eqref{eq:ch2.outgoing-idempotency} and resolvent
selection Eq.~\eqref{eq:ch2.outgoing-resolvent-selection} are consistent with
exterior radiative uniqueness (Subsection~\ref{sec:ch243}):
\begin{equation}
\label{eq:ch2.outgoing-resolvent-unique}
\mathcal{O}_{\mathrm{out}}\circ\mathcal{L}_\omega^{-1}
=
\mathcal{L}_\omega^{-1}\big|_{\mathrm{phys}}
\quad\text{is unique on}\ \mathcal{Y}_\omega,
\end{equation}
because any two selections differing by a homogeneous outgoing mode would violate
Eq.~\eqref{eq:ch2.exterior-uniqueness-theorem} \cite{stratton1941}.
Equation~\eqref{eq:ch2.outgoing-resolvent-unique} closes the excitation-to-flux sequence of
Figure~\ref{fig:ch2.outgoing-selector-measurement chain} on exterior domains: once data are
fixed, \(\mathcal{O}_{\mathrm{out}}\), \(\mathcal{S}_\omega\), and
\(\Pi_{\mathrm{rad}}\) act on a single field orbit in \(\mathcal{F}_{\mathrm{SF}}\)
\cite{stratton1941}.
On bounded domains with passive \(\mathcal{B}_\Gamma\), uniqueness follows from
Eq.~\eqref{eq:ch2.lax-milgram-wellposed} and
the passive uniqueness argument (Subsection~\ref{sec:ch01.1.3}) after quotienting
Eq.~\eqref{eq:ch2.coercivity-quotient} \cite{stratton1941}. At cavity eigenfrequencies,
\begin{equation}
\label{eq:ch2.exterior-vs-cavity-uniqueness}
\text{bounded PEC cavity:}\ \ker\mathcal{L}_\omega\neq\{\mathbf{0}\}
\quad\text{vs.}\quad
\text{exterior outgoing:}\ \ker\mathcal{L}_\omega|_{\mathcal{F}_{\mathrm{SF}}}=\{\mathbf{0}\},
\end{equation}
so Eq.~\eqref{eq:ch2.exterior-uniqueness-theorem} does not apply to closed resonators
where Eq.~\eqref{eq:ch2.cavity-resonance-failure} holds \cite{harrington2001}.
Exterior uniqueness is therefore a property of \(\mathcal{F}_{\mathrm{SF}}\) with
Sommerfeld completion, not of arbitrary \(\mathcal{B}_\Gamma\) from
Eq.~\eqref{eq:ch1.outer-boundary-classes} \cite{stratton1941}. A finite scatterer
embedded in free space may support Mie resonances in the near field without breaking
exterior uniqueness: those modes are driven by \(\Jimp\) and do not appear as free
homogeneous outgoing states \cite{jackson1999}.
Laboratory errors and non-applicability. Exterior radiative uniqueness (Subsection~\ref{sec:ch243}) requires linear passive media,
fixed \(\omega\) away from pathological embedded spectra, and exact outgoing radiation condition
rather than approximate ABC alone \cite{stratton1941}. First-order
absorbing boundaries that do not enforce Eq.~\eqref{eq:ch2.sommerfeld-field-condition}
at large \(R\) can admit spurious incoming content; uniqueness is then a property of
the \emph{augmented} truncated problem, not of the physical exterior
\cite{harrington2001}. Anisotropic, moving, or nonlinear media require modified flux
functionals; time-varying \(\mathcal{B}_\Gamma(t)\) breaks the harmonic reduction of
Subsection~\ref{sec:ch223} \cite{stratton1941}. Uniqueness of \(\mathcal{S}_\omega\)
does not imply uniqueness of \(\Jimp\) in \(\mathcal{R}_\omega\)---that is
Subsection~\ref{sec:ch244} \cite{harrington2001}. Regime qualification for transport
remains Eq.~\eqref{eq:ch2.regime-operator-screen}, independent of field uniqueness
\cite{stratton1941}.
As a worked example, truncating at finite \(R\) with a first-order ABC may return two numerically distinct
fields that agree on \(\Gamma_R\) but differ in spurious incoming content reflected
between \(\Omega_s\) and the ABC \cite{harrington2001}. Those differences are not
counterexamples to exterior radiative uniqueness (Subsection~\ref{sec:ch243}) on
\(\Omega_{\mathrm{ext}}\); they signal that the truncated BVP lacks
Eq.~\eqref{eq:ch2.sommerfeld-outgoing-class}. Silver--Müller surrogates
Eq.~\eqref{eq:ch2.silver-muller-trace-approx} improve the completion but still require
\(R\gg\lambda_0\) for uniqueness in practice \cite{stratton1941}.
